\documentclass[12pt]{article}
\usepackage[a4paper,margin=1in]{geometry}
\usepackage[utf8]{inputenc}
\usepackage[T1]{fontenc}
\usepackage{amsmath,amssymb}
\usepackage{textcomp}
\usepackage{parskip}
\begin{document}
\begin{center}
{\Large \textbf{2011 STEP II}}
\end{center}

\subsection*{Page 1}
\begin{flushleft}
\bigskip\textbf{Section A: Pure Mathematics}\par
\bigskip\textbf{Question 1.} (i) Sketch the curve y = \sqrt1 x + \sqrt3 + x .\par
-\par
Use your sketch to show that only one real value of x satisfies\par
\[\sqrt\{1\} x + \sqrt\{3\} + x = x + 1 ,\]
-\par
and give this value.\par
\textbf{(ii)} Determine graphically the number of real values of x that satisfy\par
\[2\sqrt\{1\} x = \sqrt\{3\} + x + \sqrt\{3\} x .\]
- -\par
Solve this equation.\par
2 Write down the cubes of the integers 1, 2, . . . , 10 .\par
\[The positive integers x, y and z, where x < y, satisfy\]
3 3 3\par
\[x + y = kz , ( )\]
∗\par
where k is a given positive integer.\par
\textbf{(i)} In the case x + y = k, show that\par
3 2 2\par
\[z = k 3kx + 3x .\]
-\par
\[Deduce that (4z3 k^2)/3 is a perfect square and that 1 k^2 6 z^3 < k^2 .\]
- 4\par
\[Use these results to find a solution of ( ) when k = 20.\]
∗\par
\textbf{(ii)} By considering the case x + y = z2, find two solutions of ( ) when k = 19.\par
∗\par
\[9470 Jun11\]
\end{flushleft}
\subsection*{Page 2}
\begin{flushleft}
\bigskip\textbf{Question 3.} In this question, you may assume without proof that any function f for which f 0 (x) > 0 is\par
\[increasing; that is, f(x^2) > f(x^1) if x^2 > x^1 .\]
\textbf{(i)} (a) Let f(x) = sin x x cos x. Show that f(x) is increasing for 0 6 x 6 1 \pi and deduce\par
\[that f(x) > 0 for - 0 6 x 6 1 \pi . 2\]
d\par
\[(b) Given that (arcsin x) > 1 for 0 6 x < 1, show that\]
dx\par
\[arcsin x > x (0 6 x < 1).\]
\[(c) Let g(x) = x cosec x for 0 < x < \pi. Show that g is increasing and deduce that\]
(arcsin x) x\par
-1\par
\[> x cosec x (0 < x < 1).\]
d\par
\textbf{(ii)} Given that (arctan x) 6 1 for x > 0, show by considering the function x -1 tan x that\par
dx\par
\[(tan x)(arctan x) > x 2 (0 < x < 1 \pi).\]
\textasciicircum{}\circ \textasciicircum{}\circ\par
\bigskip\textbf{Question 4.} (i) Find all the values of θ, in the range 0 < θ < 180 , for which cos θ = sin 4θ. Hence\par
show that\par
\[sin 18 ^\circ = \sqrt\{5\} 1 .\]
4 -\par
\gamma \delta\par
\textbf{(ii)} Given that\par
2 2\par
\[4 sin x + 1 = 4 sin 2x ,\]
find all possible values of sin x , giving your answers in the form p + q\sqrt5 where p and q\par
are rational numbers.\par
\textasciicircum{}\circ \textasciicircum{}\circ\par
\textbf{(iii)} Hence find two values of α with 0 < α < 90 for which\par
2 2 2\par
\[sin 3α + sin 5α = sin 6α .\]
\[9470 Jun11 [Turn over\]
\end{flushleft}
\subsection*{Page 3}
\begin{flushleft}
\bigskip\textbf{Question 5.} The points A and B have position vectors a and b with respect to an origin O, and O, A and B\par
are non-collinear. The point C, with position vector c, is the reflection of B in the line through\par
O and A. Show that c can be written in the form\par
\[c = λa b\]
-\par
2 a.b\par
\[where λ = .\]
a.a\par
The point D, with position vector d, is the reflection of C in the line through O and B. Show\par
that d can be written in the form\par
\[d = µb λa\]
-\par
for some scalar µ to be determined.\par
Given that A, B and D are collinear, find the relationship between λ and µ. In the case\par
\[λ = 1 , determine the cosine of ∠AOB and describe the relative positions of A, B and D.\]
- 2\par
\bigskip\textbf{Question 6.} For any given function f, let\par
0 2 n\par
\[I = [f (x)] [f(x)] dx , ( )\]
∗\par
\[\int\]
00 0\par
where n is a positive integer. Show that, if f(x) satisfies f (x) = kf(x)f (x) for some constant\par
k, then ( ) can be integrated to obtain an expression for I in terms of f(x), f (x), k and n.\par
∗\par
\textbf{(i)} Verify your result in the case f(x) = tan x . Hence find\par
sin x\par
dx .\par
\[cos8 x\]
\[\int\]
\textbf{(ii)} Find\par
2 6\par
sec x (sec x + tan x) dx .\par
\[\int\]
\[9470 Jun11\]
\end{flushleft}
\subsection*{Page 4}
\begin{flushleft}
\bigskip\textbf{Question 7.} The two sequences a0, a1, a2, . . . and b0, b1, b2, . . . have general terms\par
n n n n\par
\[an = λ + µ and bn = λ µ ,\]
-\par
\[respectively, where λ = 1 + \sqrt\{2\} and µ = 1 \sqrt\{2\} .\]
-\par
n n\par
\textbf{(i)} Show that br =\par
-\par
\[\sqrt\{2\} +\]
\[\sqrt\{2\}\]
a\par
n+1\par
, and give a corresponding result for ar .\par
\[r=0 r=0\]
X X\par
\textbf{(ii)} Show that, if n is odd,\par
2n m\par
1 2\par
a\par
r\par
\[=\]
bn+1 ,\par
\[m=0 r=0 !\]
X X\par
and give a corresponding result when n is even.\par
\textbf{(iii)} Show that, if n is even,\par
n 2 n\par
ar a\par
2r+1\par
\[= 2 ,\]
-\par
\[r=0 ! r=0\]
X X\par
and give a corresponding result when n is odd.\par
\[9470 Jun11 [Turn over\]
\end{flushleft}
\subsection*{Page 5}
\begin{flushleft}
\bigskip\textbf{Question 8.} The end A of an inextensible string AB of length \pi is attached to a point on the circumference\par
of a fixed circle of unit radius and centre O. Initially the string is straight and tangent to the\par
circle. The string is then wrapped round the circle until the end B comes into contact with\par
the circle. The string remains taut during the motion, so that a section of the string is in\par
contact with the circumference and the remaining section is straight.\par
Taking O to be the origin of cartesian coordinates with A at ( 1, 0) and B initially at ( 1, \pi),\par
- -\par
show that the curve described by B is given parametrically by\par
\[x = cos t + t sin t , y = sin t t cos t ,\]
-\par
where t is the angle shown in the diagram.\par
B\par
y\par
t\par
A\par
x\par
O\par
\[Find the value, t^0, of t for which x takes its maximum value on the curve, and sketch the\]
curve.\par
dx\par
Use the area integral y dt to find the area between the curve and the x axis for \pi > t > t0.\par
dt\par
\[\int\]
Find the area swept out by the string (that is, the area between the curve described by B and\par
the semicircle shown in the diagram).\par
\[9470 Jun11\]
\end{flushleft}
\subsection*{Page 6}
\begin{flushleft}
\bigskip\textbf{Section B: Mechanics}\par
\bigskip\textbf{Question 9.} Two particles, A of mass 2m and B of mass m, are moving towards each other in a straight\par
line on a smooth horizontal plane, with speeds 2u and u respectively. They collide directly.\par
Given that the coefficient of restitution between the particles is e, where 0 < e 6 1, determine\par
the speeds of the particles after the collision.\par
After the collision, B collides directly with a smooth vertical wall, rebounding and then\par
colliding directly with A for a second time. The coefficient of restitution between B and the\par
wall is f , where 0 < f 6 1. Show that the velocity of B after its second collision with A is\par
2 2 1 2\par
(1 e )u (1 4e )f u\par
3 - - 3 -\par
towards the wall and that B moves towards (not away from) the wall for all values of e and f .\par
\bigskip\textbf{Question 10.} A particle is projected from a point on a horizontal plane, at speed u and at an angle θ above\par
the horizontal. Let H be the maximum height of the particle above the plane. Derive an\par
expression for H in terms of u, g and θ.\par
A particle P is projected from a point O on a smooth horizontal plane, at speed u and at an\par
angle θ above the horizontal. At the same instant, a second particle R is projected horizontally\par
from O in such a way that R is vertically below P in the ensuing motion. A light inextensible\par
string of length H connects P and R. Show that the time that elapses before the string\par
becomes taut is\par
\[(\sqrt\{2\} 1) H/g .\]
-\par
p\par
When the string becomes taut, R leaves the plane, the string remaining taut. Given that P\par
and R have equal masses, determine the total horizontal distance, D, travelled by R from the\par
moment its motion begins to the moment it lands on the plane again, giving your answer in\par
terms of u, g and θ.\par
\[Given that D = H, find the value of tan θ.\]
\[9470 Jun11 [Turn over\]
\end{flushleft}
\subsection*{Page 7}
\begin{flushleft}
\bigskip\textbf{Question 11.} Three non-collinear points A, B and C lie in a horizontal ceiling. A particle P of weight W\par
is suspended from this ceiling by means of three light inextensible strings AP , BP and CP ,\par
as shown in the diagram. The point O lies vertically above P in the ceiling.\par
A C\par
O\par
B\par
P\par
\textasciicircum{}\circ \textasciicircum{}\circ\par
The angles AOB and AOC are 90 + θ and 90 + φ, respectively, where θ and φ are acute\par
\[angles such that tan θ = \sqrt\{2\} and tan φ = 1 \sqrt\{2\}.\]
\textasciicircum{}\circ \textasciicircum{}\circ \textasciicircum{}\circ\par
The strings AP , BP and CP make angles 30 , 90 θ and 60 , respectively, with the vertical,\par
-\par
and the tensions in these strings have magnitudes T , U and V respectively.\par
\textbf{(i)} Show that the unit vector in the direction P B can be written in the form\par
\[1 \sqrt\{2\} \sqrt\{2\}\]
i j + k ,\par
\[- 3 - 3 \sqrt\{3\}\]
where i , j and k are the usual mutually perpendicular unit vectors with j parallel to\par
OA and k vertically upwards.\par
\textbf{(ii)} Find expressions in vector form for the forces acting on P .\par
\textbf{(iii)} Show that U = \sqrt6V and find T , U and V in terms of W .\par
\[9470 Jun11\]
\end{flushleft}
\subsection*{Page 8}
\begin{flushleft}
\bigskip\textbf{Section C: Probability and Statistics}\par
12 Xavier and Younis are playing a match. The match consists of a series of games and each\par
game consists of three points.\par
Xavier has probability p and Younis has probability 1 p of winning the first point of any\par
-\par
game. In the second and third points of each game, the player who won the previous point\par
has probability p and the player who lost the previous point has probability 1 p of winning\par
-\par
the point. If a player wins two consecutive points in a single game, the match ends and that\par
player has won; otherwise the match continues with another game.\par
\textbf{(i)} Let w be the probability that Younis wins the match. Show that, for p = 0,\par
\[1 p^2\]
\[w = - .\]
2 p\par
-\par
1 1 1 1\par
\[Show that w > if p < , and w < if p > . Does w increase whenever p decreases?\]
2 2 2 2\par
\textbf{(ii)} If Xavier wins the match, Younis gives him \pounds 1; if Younis wins the match, Xavier gives\par
\[him \pounds k. Find the value of k for which the game is fair in the case when p = 2 .\]
\textbf{(iii)} What happens when p = 0?\par
\[9470 Jun11 [Turn over\]
\end{flushleft}
\subsection*{Page 9}
\begin{flushleft}
13 What property of a distribution is measured by its skewness?\par
\textbf{(i)} One measure of skewness, γ, is given by\par
E (X µ)3\par
\[γ = - ,\]
σ3\par
where µ and σ2 are the mean and variance of the random variable X. Show that\par
\[E(X^3) 3µσ2 µ3\]
\[γ = - - .\]
σ3\par
The continuous random variable X has probability density function f where\par
2x for 0 6 x 6 1 ,\par
\[f(x) =\]
\par
0 otherwise .\par
\par
\[2\sqrt\{2\}\]
\[Show that for this distribution γ = .\]
- 5\par
\textbf{(ii)} The decile skewness, D, of a distribution is defined by\par
F\par
-1( 9\par
) 2F\par
-1( 1\par
) + F\par
-1( 1\par
)\par
\[D = 10 - 2 10 ,\]
F-1( 9 ) F-1( 1 )\par
10 - 10\par
where F\par
-1\par
is the inverse of the cumulative distribution function. Show that, for the\par
\[above distribution, D = 2 \sqrt\{5\} .\]
-\par
The Pearson skewness, P , of a distribution is defined by\par
3(µ M )\par
\[P = - ,\]
σ\par
\[where M is the median. Find P for the above distribution and show that D > P > γ .\]
\[9470 Jun11\]
\end{flushleft}
\end{document}
